\documentclass[aspectratio=169]{beamer}
\usetheme{metropolis}
\usepackage{appendixnumberbeamer}
\usepackage{booktabs, hyperref}

\input{mgmt675-style}

\usepackage{tikz}
\usetikzlibrary{shapes.geometric, arrows.meta, positioning, calc}

% Title info
\subtitle{MGMT 675: Generative AI for Finance}
\title{Module 3: Getting Started with Claude Code}
\author{Kerry Back}
\date{}

\begin{document}

\maketitle

% ========================================
% SECTION 1: CLAUDE DESKTOP OVERVIEW
% ========================================
\section{Claude Desktop Overview}

\begin{frame}{Three Tabs: Chat, Cowork, Code}
\begin{columns}[T]
\begin{column}{0.32\textwidth}
\begin{shadedbox}[title=\textbf{Chat (Analysis)}]
\begin{itemize}\footnotesize
  \item Sandboxed Python in the browser
  \item No local file access
  \item Must upload data manually
\end{itemize}
\end{shadedbox}
\end{column}
\begin{column}{0.32\textwidth}
\begin{shadedbox}[title=\textbf{Cowork}]
\begin{itemize}\footnotesize
  \item Runs in a local VM
  \item Files synced to/from VM
  \item \alert{No internet access}
\end{itemize}
\end{shadedbox}
\end{column}
\begin{column}{0.32\textwidth}
\begin{shadedbox}[title=\textbf{Code}]
\begin{itemize}\footnotesize
  \item Runs on \alert{your machine}
  \item Full local file access
  \item Full internet access
\end{itemize}
\end{shadedbox}
\end{column}
\end{columns}
\vspace{0.3cm}
\begin{center}
\alert{Code mode} gives Claude direct access to your machine --- no sandbox, no VM, no upload step.
\end{center}
\end{frame}

\begin{frame}{Which Tool When?}
\begin{itemize}
  \item \textbf{Chat} for questions and explanations (e.g., ``Explain the WACC formula'')
  \item \textbf{Artifacts} for interactive tools (e.g., a DCF calculator with sliders)
  \item \textbf{Code} for anything requiring files, APIs, or local execution (e.g., fetch live data, run regressions, create Excel)
  \item \textbf{Cowork} for complex multi-file tasks with no internet needed (e.g., organize 50 PDF 10-Ks)
\end{itemize}
\vspace{0.3cm}
\textbf{Rule of thumb:} Start with Chat. If you need a visual, use Artifacts. If you need file I/O, use Code. Save Cowork for complex multi-file tasks.
\end{frame}

% ========================================
% SECTION 2: SETTING UP CLAUDE CODE
% ========================================
\section{Setting Up Claude Code}

\begin{frame}{Running the Installer}
The class installer sets up everything you need in one step:
\vspace{0.3cm}
\begin{itemize}
  \item \textbf{Python} --- language Claude Code uses to execute analysis
  \item \textbf{Git} and \textbf{GitHub CLI} --- version control and collaboration
  \item \textbf{Claude Code} --- the agent engine that runs on your machine
  \item \textbf{Skills} --- pre-packaged instructions for Excel, Word, PowerPoint, and more
\end{itemize}
\vspace{0.3cm}
\alert{Follow the class setup guide for your platform (Mac or Windows).}
\end{frame}

\begin{frame}{Your First Session}
\begin{columns}[T]
\begin{column}{0.5\textwidth}
\begin{shadedbox}[title=\textbf{Getting Started}]
\begin{enumerate}\small
  \item Open Claude Desktop $\rightarrow$ \textbf{Code} tab
  \item Select a working folder for your project
  \item Choose a model: \textbf{Sonnet} (fast) or \textbf{Opus} (thorough)
  \item Set permission mode to \textbf{auto-accept} (recommended for exercises)
\end{enumerate}
\end{shadedbox}
\end{column}
\begin{column}{0.5\textwidth}
\begin{shadedbox}[title=\textbf{Try It}]
\begin{itemize}\small
  \item ``Create a CSV file with five rows of sample sales data --- store name, date, product, quantity, revenue --- then load it into a pandas DataFrame and plot revenue by store.''
  \item Claude writes the code, runs it, and shows you the result
\end{itemize}
\end{shadedbox}
\end{column}
\end{columns}
\end{frame}

\begin{frame}{Terminal Alternative}
The Code tab in Claude Desktop and the terminal command \texttt{claude} use the \alert{same engine}.
\vspace{0.3cm}
\begin{columns}[T]
\begin{column}{0.5\textwidth}
\begin{shadedbox}[title=\textbf{Claude Desktop (Code tab)}]
\begin{itemize}\small
  \item Visual interface with file browser
  \item Click to select working folder
  \item Good for getting started
\end{itemize}
\end{shadedbox}
\end{column}
\begin{column}{0.5\textwidth}
\begin{shadedbox}[title=\textbf{Terminal (\texttt{claude})}]
\begin{itemize}\small
  \item Type \texttt{claude} in any terminal
  \item Same skills, same config
  \item Preferred by power users
\end{itemize}
\end{shadedbox}
\end{column}
\end{columns}
\vspace{0.3cm}
Both approaches share the same skills, MCP servers, and configuration files.
\end{frame}

% ========================================
% SECTION 3: CLAUDE.MD AND SKILLS
% ========================================
\section{CLAUDE.md and Skills}

\begin{frame}{Why Skills Matter for Finance}
When you build DCF models and portfolio optimizations in ad-hoc conversations, what happens when you start a new conversation?
\vspace{0.3cm}
\begin{columns}[T]
\begin{column}{0.5\textwidth}
\begin{shadedbox}[title=\textbf{The Problem}]
\begin{itemize}\small
  \item One conversation uses 10\% WACC, another uses 8\%
  \item One assumes 3\% terminal growth, another 2.5\%
  \item Portfolio optimization: 5-year history vs.\ 3-year
\end{itemize}
\end{shadedbox}
\end{column}
\begin{column}{0.5\textwidth}
\begin{shadedbox}[title=\textbf{The Fix: Skills}]
\begin{itemize}\small
  \item A DCF skill specifies: which assumptions, how to estimate them, what output format
  \item A portfolio skill specifies: data source, estimation window, constraints
  \item \alert{Every conversation follows the same playbook}
\end{itemize}
\end{shadedbox}
\end{column}
\end{columns}
\vspace{0.3cm}
\alert{No audit trail, no reproducibility --- unless you encode the methodology in a skill.}
\end{frame}

\begin{frame}{Finance Skill Examples}
\begin{columns}[T]
\begin{column}{0.5\textwidth}
\begin{shadedbox}[title=\textbf{DCF Skill}]
\begin{itemize}\small
  \item Sales growth from trailing 5-year CAGR
  \item Sector median EV/EBITDA for exit multiple
  \item Two-way sensitivity table (WACC vs.\ growth)
\end{itemize}
\end{shadedbox}
\end{column}
\begin{column}{0.5\textwidth}
\begin{shadedbox}[title=\textbf{Investment Memo Skill}]
\begin{itemize}\small
  \item Standard structure: exec summary, financials, risks, recommendation
  \item Required data points and formatting
  \item Every memo follows the same template
\end{itemize}
\end{shadedbox}
\end{column}
\end{columns}
\vspace{0.3cm}
\alert{Every analyst produces comparable, auditable outputs.}
\end{frame}

\begin{frame}{CLAUDE.md: Global vs.\ Per-Project}
Claude reads \texttt{CLAUDE.md} files automatically at the start of every session --- no need to repeat instructions.
\vspace{0.3cm}
\begin{columns}[T]
\begin{column}{0.5\textwidth}
\begin{shadedbox}[title=\textbf{Global (\textasciitilde/.claude/CLAUDE.md)}]
\begin{itemize}\small
  \item Applies to \alert{every} project
  \item Output style preferences (chart colors, file formats)
  \item Preferred libraries (e.g., ``use pandas, not polars'')
  \item Personal conventions
\end{itemize}
\end{shadedbox}
\end{column}
\begin{column}{0.5\textwidth}
\begin{shadedbox}[title=\textbf{Per-Project (project root)}]
\begin{itemize}\small
  \item Applies only to \alert{this} folder
  \item Data context (column names, units, sources)
  \item Domain guidance (industry terms, formulas)
  \item Project-specific workflows
\end{itemize}
\end{shadedbox}
\end{column}
\end{columns}
\vspace{0.3cm}
Both are plain text files.  Claude reads the global file first, then the project file.  Project instructions can override global ones.
\end{frame}

\begin{frame}{Pre-Installed Skills: Document Creation}
\textbf{Skills} are packaged instructions that teach Claude Code a specific workflow.  The class installer pre-configured four.  Invoke any skill by typing \texttt{/} followed by its name.
\vspace{0.3cm}
\begin{columns}[T]
\begin{column}{0.32\textwidth}
\begin{shadedbox}[title=\textbf{/xlsx}]
\begin{itemize}\footnotesize
  \item Create and edit Excel workbooks
  \item Uses \alert{live formulas}, not hardcoded values
  \item Supports charts, formatting, multiple sheets
\end{itemize}
\end{shadedbox}
\end{column}
\begin{column}{0.32\textwidth}
\begin{shadedbox}[title=\textbf{/docx}]
\begin{itemize}\footnotesize
  \item Create and edit Word documents
  \item Professional formatting, headers, tables
  \item Supports tracked changes and comments
\end{itemize}
\end{shadedbox}
\end{column}
\begin{column}{0.32\textwidth}
\begin{shadedbox}[title=\textbf{/pptx}]
\begin{itemize}\footnotesize
  \item Create and edit PowerPoint decks
  \item Slide layouts, charts, speaker notes
  \item Can modify existing presentations
\end{itemize}
\end{shadedbox}
\end{column}
\end{columns}
\vspace{0.3cm}
\textbf{Example:} \textit{``/xlsx Create a loan amortization table for a \$200K mortgage at 6.5\% with monthly payments and a chart of principal vs.\ interest.''}
\end{frame}

\begin{frame}{Pre-Installed Skills: Skill Creator}
The \texttt{/skill-creator} skill is a guide for \alert{building your own skills}.
\vspace{0.3cm}
\begin{columns}[T]
\begin{column}{0.5\textwidth}
\begin{shadedbox}[title=\textbf{What It Does}]
\begin{itemize}\small
  \item Walks you through designing a new skill step by step
  \item Generates the skill definition file
  \item Installs the skill so it appears as a \texttt{/} command
\end{itemize}
\end{shadedbox}
\end{column}
\begin{column}{0.5\textwidth}
\begin{shadedbox}[title=\textbf{When to Use It}]
\begin{itemize}\small
  \item You repeat the same workflow across sessions
  \item You want Claude to follow a specific multi-step process
  \item You want to share a workflow with classmates
\end{itemize}
\end{shadedbox}
\end{column}
\end{columns}
\vspace{0.3cm}
We will use \texttt{/skill-creator} in the exercise below to build a personal task manager skill.
\end{frame}

\begin{frame}{Exercise: Create a Tasks Skill}
Use \texttt{/skill-creator} to build a \textbf{tasks} skill that manages personal and school to-do lists.
\vspace{0.3cm}
\begin{columns}[T]
\begin{column}{0.5\textwidth}
\begin{shadedbox}[title=\textbf{What the Skill Should Do}]
\begin{itemize}\small
  \item Store tasks in markdown files (e.g., \texttt{personal.md}, \texttt{school.md})
  \item Add new tasks with a category and due date
  \item List pending tasks when asked
  \item Mark tasks as complete
\end{itemize}
\end{shadedbox}
\end{column}
\begin{column}{0.5\textwidth}
\begin{shadedbox}[title=\textbf{Steps}]
\begin{enumerate}\small
  \item Type \texttt{/skill-creator} and describe the task manager you want
  \item Let Claude generate and install the skill
  \item Test it: ``/tasks Add: finish DCF exercise, due Tuesday''
  \item Test it: ``/tasks What's pending?''
\end{enumerate}
\end{shadedbox}
\end{column}
\end{columns}
\vspace{0.3cm}
This exercise teaches skill creation.  The same pattern applies to building skills for data retrieval, report generation, or any repeated workflow.
\end{frame}

% ========================================
% SECTION 4: WORKING WITH MULTIPLE FILES
% ========================================
\section{Working with Multiple Files}

\begin{frame}{Merging Data from Different Sources}
Real financial data rarely lives in a single file or format.  Claude Code can read CSV, Excel, PDF, and Word files, then \alert{combine them into a unified dataset}.
\vspace{0.3cm}
\begin{columns}[T]
\begin{column}{0.5\textwidth}
\begin{shadedbox}[title=\textbf{Common Scenarios}]
\begin{itemize}\small
  \item Holdings from multiple brokerages (xlsx + csv)
  \item Loan tape (CSV) + appraisal report (PDF) + policy memo (DOCX)
  \item Quarterly reports from divisions with different column names
\end{itemize}
\end{shadedbox}
\end{column}
\begin{column}{0.5\textwidth}
\begin{shadedbox}[title=\textbf{What Claude Handles}]
\begin{itemize}\small
  \item Reads each file format natively
  \item Reconciles column names and data types
  \item Flags mismatches and missing values
  \item Outputs a clean, unified table
\end{itemize}
\end{shadedbox}
\end{column}
\end{columns}
\vspace{0.3cm}
\alert{You describe what should be merged and how.  Claude handles the parsing and code.}
\end{frame}

\begin{frame}{Filtering and Aggregating}
Once data is combined, the next step is usually filtering, sorting, or grouping.
\vspace{0.3cm}
\begin{columns}[T]
\begin{column}{0.5\textwidth}
\begin{shadedbox}[title=\textbf{Filtering}]
\begin{itemize}\small
  \item ``Keep only loans above \$1M''
  \item ``Filter to stocks with P/E > 0''
  \item ``Exclude rows with missing revenue''
\end{itemize}
\end{shadedbox}
\end{column}
\begin{column}{0.5\textwidth}
\begin{shadedbox}[title=\textbf{Aggregating}]
\begin{itemize}\small
  \item ``Average revenue by sector and quarter''
  \item ``Sum holdings by asset class''
  \item ``Rank funds by Sharpe ratio''
\end{itemize}
\end{shadedbox}
\end{column}
\end{columns}
\vspace{0.3cm}
State your intent in plain English.  Claude writes the pandas code, runs it, and shows you the result.  Always \alert{check the filter logic} --- this is where silent errors happen most often (Module~5).
\end{frame}

\begin{frame}{Exercise: Portfolio Consolidation}
Download \href{https://kerryback.com/mgmt675/files/portfolio-consolidation.zip}{portfolio-consolidation.zip} and extract the files into a folder.  You will have holdings data from three brokerages in different formats (two Excel workbooks and one CSV).
\vspace{0.3cm}
In Claude Code, ask Claude to:
\begin{enumerate}
  \item Merge all equity holdings into a single consolidated view with standardized columns
  \item Resolve inconsistencies between ticker symbols and full company names
  \item Sum shares for positions that appear in multiple accounts
  \item Produce a summary showing total portfolio value, top five holdings by market value, and asset allocation between equities and bonds
\end{enumerate}
\vspace{0.2cm}
Save the result as \texttt{consolidated.xlsx}.
\end{frame}

\begin{frame}{Exercise: Aggregating Diverse File Types}
Download \href{https://kerryback.com/mgmt675/files/loans.zip}{loans.zip} and extract the files into a folder.  You will have a loan tape (CSV), a collateral appraisal report (PDF), and a policy exceptions memo (DOCX).
\vspace{0.3cm}
Ask Claude Code to read all three files and produce a single Excel workbook that:
\begin{itemize}
  \item Lists every loan with its terms from the CSV
  \item Adds a column with the appraised collateral value from the PDF
  \item Flags loans that violate the policy limits described in the memo
\end{itemize}
\vspace{0.3cm}
\alert{Notice that Claude reads CSV, PDF, and DOCX natively} --- no file-parsing code, no manual extraction, no copy-paste.
\end{frame}

\begin{frame}{Exercise: Reconciling Column Names}
Download \href{https://kerryback.com/mgmt675/files/aggregation.zip}{aggregation.zip} and extract the workbooks into a folder.
\vspace{0.3cm}
Each Excel file contains a data table.  Column names vary across files and some columns are missing.  Ask Claude Code to:
\begin{enumerate}
  \item Summarize the columns present, row count, and basic statistics for each file
  \item Combine all files into a single table, reconciling the varying column names
  \item Save the result as \texttt{combined.xlsx}
\end{enumerate}
\vspace{0.3cm}
This is the same ``divide and conquer'' pattern used in professional data pipelines --- read each source independently, then unify.
\end{frame}

% ========================================
% SECTION 5: EFFECTIVE PROMPTING
% ========================================
\section{Effective Prompting}

\begin{frame}{The Plan-Execute-Evaluate Cycle}
\begin{center}
\begin{tikzpicture}[scale=1.2]
  % Nodes
  \node[draw, rounded corners, fill=accentblue!20, minimum width=2.5cm, minimum height=1cm] (plan) at (0,2) {\textbf{Plan}};
  \node[draw, rounded corners, fill=accentblue!20, minimum width=2.5cm, minimum height=1cm] (execute) at (4,2) {\textbf{Execute}};
  \node[draw, rounded corners, fill=accentblue!20, minimum width=2.5cm, minimum height=1cm] (evaluate) at (2,0) {\textbf{Evaluate}};

  % Arrows
  \draw[->, thick, accentblue] (plan) -- (execute);
  \draw[->, thick, accentblue] (execute) -- (evaluate);
  \draw[->, thick, accentblue] (evaluate) -- (plan);
\end{tikzpicture}
\end{center}

\begin{itemize}
  \item \textbf{Plan:} Define the task, gather requirements, outline approach
  \item \textbf{Execute:} Let AI do the work with your guidance
  \item \textbf{Evaluate:} Check results, identify gaps, refine---then repeat
\end{itemize}
\end{frame}

\begin{frame}{Cross-Evaluation}
\begin{itemize}
  \item Start a \alert{new conversation} to evaluate the plan and output---even with the same model, a fresh context can catch errors
  \item Ask the new conversation to review, critique, or verify; different phrasing may reveal blind spots
  \item Consider using a \textit{different model} for additional perspective
\end{itemize}

\vspace{0.3cm}
\textbf{Example:} \textit{``I received this analysis from another session. Can you review it for errors or questionable assumptions?''}

\vspace{0.3cm}
This technique becomes critical in Module 5 (Verification and Governance).
\end{frame}

\end{document}
